% Source: Solutions/Sahil-PS1-Q6-10.pdf, p. 1 (typed). The opening convention paragraph is
% the preface of Sahil's whole submission (Q6-10).
\submittedby{Sahil}

Throughout, a topology on a set $X$ is a collection $\T\subseteq\Pow(X)$ such that
(1) $\varnothing,X\in\T$, (2) arbitrary unions of members of $\T$ belong to $\T$, and
(3) finite intersections of members of $\T$ belong to $\T$.

Let
\[ A_n=\{1,2,\dots,n\},\qquad B_n=\N\setminus A_n\quad(n\ge1). \]
Then the two collections in the question can be written as
\[ \T_1=\{\varnothing,\N\}\cup\{A_n:n\ge1\},\qquad \T_2=\{\varnothing,\N\}\cup\{B_n:n\ge1\}. \]
We verify the topology axioms separately.

\textbf{(a) The collection $\T_1$ is a topology.} Clearly, $\varnothing,\N\in\T_1$. For finite
intersections, observe that
\[ A_m\cap A_n=A_{\min\{m,n\}}. \]
Intersections with $\varnothing$ or $\N$ also belong to $\T_1$. Thus $\T_1$ is closed under
finite intersections.

Now let $\{U_i\}_{i\in I}$ be any family of members of $\T_1$. If every $U_i$ is empty, then
$\bigcup_{i\in I}U_i=\varnothing$. If at least one $U_i=\N$, the union is $\N$. In the remaining
case, the nonempty sets in the family are of the form $A_n$. If their indices are unbounded in
$\N$, then every positive integer belongs to some $A_n$, so their union is $\N$. If their
indices are bounded, they have a largest value $m$ (a nonempty bounded subset of $\N$ has a
maximum), and the union is $A_m$. In every case the union belongs to $\T_1$. Hence $\T_1$ is a
topology on $\N$.

\textbf{(b) The collection $\T_2$ is a topology.} Clearly, $\varnothing,\N\in\T_2$. For finite
intersections, De Morgan's law gives
\[ B_m\cap B_n=(\N\setminus A_m)\cap(\N\setminus A_n)=\N\setminus(A_m\cup A_n)=B_{\max\{m,n\}}. \]
Intersections involving $\varnothing$ or $\N$ also remain in $\T_2$.

Consider an arbitrary family $\{U_i\}_{i\in I}\subseteq\T_2$. If one member is $\N$, its union
is $\N$. If every member is empty, its union is empty. Otherwise, discard the empty members and
write the remaining members as $B_{n_i}$. The set of indices
$\{n_i : i\in I,\ U_i\neq\varnothing\}$ is a nonempty subset of $\N$, so it has a least element,
say $m$. Since the sets $B_n$ increase as $n$ decreases, we have
\[ \bigcup_{i\in I}B_{n_i}=B_m. \]
Indeed, $B_m$ is one of the sets in the union and contains every $B_{n_i}$ because $m\le n_i$.
Thus the union belongs to $\T_2$. Therefore $\T_2$ is also a topology on $\N$.

\textbf{Conclusion:} Both $\T_1$ and $\T_2$ define topologies on $\N$.
